\chapter{The Mirror: The West Runs the Playbook Against Itself}
\label{ch:mirror}

The chapters of Part~II examined the elements of AI one at a time and found, in each, a recognizable move from the playbook of Chapter~\ref{ch:playbook} being executed by a Chinese actor. This chapter reports the finding that changes what those elements mean. Over the course of 2026 the largest Western firms in the sector converged on the same method---not in imitation, and in most cases without any reference to China at all, but by the same economic logic arriving at the same instrument. Each of them now commoditizes, deliberately and at its own near-term expense, a layer it does not own, in order to defend the layer it does. The accelerator vendor gives away models. The model laboratories design accelerators. The hyperscalers build the fabric out of merchant parts and publish the specifications. A launch company is building a fab. The strategy the preceding chapters attribute to Beijing is no longer Chinese. What remains asymmetric is not the method but the scope---and, as Section~\ref{sec:mirror-asymmetry} argues, the asymmetry of scope is worth more than the symmetry of method takes away.

Two things follow immediately, and they point in opposite directions. The first is that Proposition~P3---the destruction of the rent that finances the Western frontier---acquires a second and entirely independent mechanism, one that does not require Chinese price pressure to reach Western margins at all, because a layer can be commoditized by the firm sitting above it and every layer in this stack has a firm sitting above it. The second is that a West which commoditizes its own layers competently captures the surplus of a cheap-intelligence world as surely as a China that does, and the book must not confuse falling prices with a change of ownership. Both readings are printed here, and Section~\ref{sec:mirror-license} states which claims the chapter licenses and which it forbids.

\section{The move, stated generally}
\label{sec:mirror-move}

The theory chapter established the result this chapter depends on: openness relocates rent, it does not abolish it (Section~\ref{sec:io}). A firm that gives away a good it produces is not destroying value; it is moving the point at which value is captured to a layer it holds more securely. That result was stated there as an account of Chinese strategy, but nothing in it is national. It is a statement about complements. A firm that owns layer~$A$ and faces a rent-taker at the adjacent layer~$B$ has an option: drive the price of $B$ to marginal cost, and the surplus that $B$ was extracting flows to $A$ instead. The move costs the firm nothing it was earning, transfers the loss to a competitor it does not own, and---this is the part that makes it attractive rather than merely aggressive---expands the total market, because the complement's demand rises as its own price falls.

What is new in 2026 is not the theory, which is old, but the simultaneity. Three distinct classes of Western actor executed the move at scale within a single year and a fourth announced it; one of the four stated the rationale on the record, and the others are inferred from what they built rather than from what they said. The result is a stack in which every layer is being commoditized by somebody, which raises the question this chapter is ultimately about: when everyone is giving away everything, where does the rent actually land?

\section{The accelerator vendor commoditizes the model, and underwrites the buildout}
\label{sec:mirror-nvidia}

The clearest instance is also the largest company in the sector, and it runs the move on two layers at once.

\subsection{The model layer}

NVIDIA released the Nemotron~3 family in December 2025---Nano, Super and Ultra---and released with it not only open weights but the \emph{training datasets and the reinforcement-learning libraries}~\cite{nemotron}. Nemotron~3.5 Lightning followed in August 2026 under the OpenMDW-1.1 licence, again with training data and recipes, distributed through Hugging Face, ModelScope, OpenRouter, SageMaker JumpStart and Google Cloud. The associated open pretraining collection exceeds ten trillion tokens, and NVIDIA publishes open robotics, protein-structure and autonomous-driving corpora alongside it. On the software side, TensorRT-LLM is Apache~2.0, and in December 2025 NVIDIA open-sourced the CUDA Tile intermediate representation and the cuTile compiler stack under Apache~2.0 with LLVM exceptions, describing it as the largest update to the CUDA platform since its invention~\cite{cuda-tile}.

The scale of the resulting position is measurable rather than rhetorical. In the week to 29 August 2026, Nemotron~3 Ultra was the sixth-largest model by tokens routed on the largest neutral router---on its \emph{free} endpoint---ahead of every model from every Chinese laboratory except four~\cite{openrouter-aug26} [A]. The caveat of Section~\ref{sec:price-inversion} applies here as it applies there: router volumes are a small and promotion-distorted slice of global inference, and a free endpoint is the promotion in its purest form. Read as a direction rather than a census, a chip company is a top-ten model provider by volume, which is not a description anyone would have written of this industry two years ago.

The strategic rationale is on the record, twice, and neither statement is ambiguous. Jensen Huang, on the second-quarter fiscal 2027 earnings call: ``The world will need both closed models and open models. And both closed models and open models are skyrocketing in use\ldots\ \emph{Nearly all open models run on NVIDIA}\ldots\ our position in open models is very good, because the CUDA ecosystem is literally everywhere''~\cite{nvda-open-models} [P]. And, more precisely still, NVIDIA's Director of Developer Technology at the Six Five Summit on 27~August 2026: NVIDIA builds Nemotron because ``this is going to be a multi-model world, and we see where there's the bottleneck''---arguing that durable value accrues to proprietary data rather than to model weights, which is exactly why NVIDIA can afford to release the weights \emph{and} the data~\cite{nvda-open-models} [P].

What changed in August 2026 is that NVIDIA stopped merely publishing weights and began buying the means of producing them. On 20--24~August 2026---the sources differ on the date and the difference is not resolved here---NVIDIA and Poolside disclosed a transaction of approximately \$7B: \$6B for a \emph{non-exclusive} licence to Poolside's ``Model Factory'' training platform, \$1B of equity at a \$12B pre-money valuation, and offers to 109 Poolside engineers and researchers to join NVIDIA and work on Nemotron~\cite{nvda-poolside} [P]. Poolside remains independent, its founders do not join, and the licence excludes no one else from taking one. The reported purpose is stated in terms this chapter has to take seriously: to build a leading American open-weight model capable of competing with DeepSeek and Moonshot.

That is a description of a market rather than an aspiration, and the market can be dated to the same fortnight. On 28~August 2026 Tencent open-sourced the Hy4 preview---770B total parameters, 49B active, a context window above one million tokens---the largest open-weight release of the last week of August, and three of the four open-weight releases in that window were Chinese~\cite{hunyuan-hy4} [V]. DeepSeek had meanwhile taken its first outside capital in June---roughly \$7.4B at a valuation above \$50B, with the state's Big Fund as the sole voting holder~\cite{deepseek-round} [A]; a late-August report of a \$74B valuation ahead of a 2027 onshore listing conflicts with that account and is not used here. The layer NVIDIA is buying its way into is not vacant. It is the most crowded layer in the industry and it is disproportionately Chinese, which is the condition that makes a \$6B licence for a training platform legible as something other than eccentricity.

Two days before the Tencent release, on 26~August 2026, \emph{The Information} reported that NVIDIA was closing on an acquisition of Hugging Face for approximately \$12.9B~\cite{nvda-hf} [A]. The report is a report and nothing more: neither company had confirmed it at the currency date, and \emph{Business Insider} reported that talks had produced no signed agreement and that the transaction ``could still atomize.'' This book does not assert that it happened, and the analysis that follows is written to survive its not happening. What makes it worth printing unconfirmed is the reversal it would represent. Hugging Face reportedly \emph{declined} a \$500M NVIDIA investment at a \$7B valuation in 2025, on the explicit ground that it did not want a dominant investor able to sway its decisions~\cite{nvda-hf}. A hub that refused a minority stake on independence grounds in one year and is reported in acquisition talks with the same buyer in the next is either a firm that changed its mind about independence or a firm that discovered what independence costs; the record at \datafreeze does not say which, and the possibility that the report is simply wrong remains open.

Take the three together---Nemotron, the Model Factory licence with its 109 researchers, and, if the report is right, the hub---and the description this chapter has used so far stops being adequate. A firm that publishes open weights as a loss-leader is subsidizing a complement, which is the move of Section~\ref{sec:mirror-move} and nothing more. A firm that acquires the training platform on which open models are produced, the researchers who produce them, and the distribution point through which very nearly all open weights reach their users is doing something categorically larger: it is acquiring the \emph{supply chain} of the open-model layer---the factory, the staff and the warehouse. That is a different object from a commons, and the chapter should say what it is instead. Linux is a commons in the strict sense: the licence is irrevocable, no single vendor's failure can withdraw the kernel, and the distributions that grew around it competed with each other for the surplus. An open-model layer whose training platform is licensed from one accelerator vendor, whose senior researchers are on that vendor's payroll, and whose distribution hub is that vendor's subsidiary is not a commons in that sense. It is a vendor-controlled distribution: free to use, forkable in principle, and shaped in practice by one balance sheet's judgements about what gets trained, on whose silicon, and what appears on the front page.

Both readings are available and this chapter declines to choose. On the first, that is simply how a complement gets supplied reliably. Open weights of frontier quality do not appear because a market would like them to; they appear because somebody pays for the compute, the corpus and the people, and the natural payer is a firm whose revenue rises with inference volume. The arrangement is unusual only in that the payer says out loud why it is paying. On the second, the same facts describe the concentration of an entire layer under a single owner---and the layer in question is the one this book has treated throughout as the principal solvent of proprietary advantage, which makes its ownership a different kind of fact from the ownership of a chip design. Note that the two readings assign opposite meanings to the same detail: the non-exclusivity of the Poolside licence is, on the first, proof that the platform stays available to rivals, and on the second, the feature that moves the people who know how to operate it while keeping the transaction outside merger review.

Which raises the structural question, and it is not this book's to settle. The Poolside transaction is the third of a kind. In December 2025 NVIDIA took a non-exclusive licence to Groq's inference technology together with roughly ninety per cent of Groq's staff---including its chief executive and its president---for approximately \$20B, leaving Groq to continue under new management; in 2026 it took Enfabrica's technology and team for approximately \$900M; in August 2026 it licensed Poolside's platform and hired 109 of its researchers for approximately \$7B~\cite{nvda-structured} [V]. Roughly \$27B in nine months, each transaction structured so that the target stays nominally independent and, on the reported structure, outside Hart--Scott--Rodino merger review. Senators Warren and Blumenthal put the objection in a sentence on 23~March 2026: ``by licensing its technology and hiring its most important employees, NVIDIA has effectively acquired Groq in all but name''~\cite{nvda-structured} [P]. No enforcement action had been taken by the currency date, and a letter from two senators is not a finding of law; the structures may well be lawful and are certainly not secret. The observation this chapter adds is narrower and follows from the form rather than from anyone's motive. An equity acquisition of Hugging Face \emph{would} require Hart--Scott--Rodino review, which is precisely why the report matters more than its price. Every previous transaction in the sequence was built so that the question would not be asked. This one, if it is real, asks it.

Set that beside Chapter~\ref{ch:model}. The argument for releasing weights at marginal price, on the ground that the rent will land somewhere else, is the argument this book attributed to Chinese laboratories and to the industrial-policy documents behind them. It is here being made, in public, by the incumbent whose rent the Chinese version of the move is aimed at. The two are not doing the same thing to the same end---NVIDIA is defending silicon, and China is defending an eventual position in hardware volume, domestic cloud and standards---but they are doing the same thing.

\subsection{The datacentre layer}

The second move is larger and less familiar, and it is the one that gives the first its durability. Through 2026 NVIDIA converted AI compute from a capital good its customers had to finance into an asset class that third-party capital would finance, and it did so by standing behind the residual value.

The filed positions, which supersede a good deal of looser reporting, are these. On 17~August 2026 NVIDIA disclosed multiple residual value guaranties with SB~Energy on datacentre leases at Pike County, Ohio, whose tenant is an OpenAI affiliate: approximately 4.25\,GW of IT load initially with an option on a further 3.8\,GW, a twenty-year maximum term, and an aggregate payment obligation ``cumulatively capped at \$105 billion''~\cite{nvda-ohio} [P]. The quarterly filings put maximum gross guarantee exposure at \$108.5B, supply and capacity purchase commitments at \$279B, AI-cloud revenue-share commitments at \$36B, future equity investment commitments at \$25B, datacentre leases held for third-party reassignment at \$20B, and equity securities at roughly \$95.6B~\cite{nvda-exposure} [P]. The same Ohio filing records that NVIDIA invests \$1.5B in SB~Energy and becomes the site's exclusive AI compute infrastructure provider, with the first phases effective in fiscal 2029~\cite{nvda-ohio}. Alongside these sit a residual-purchase obligation to CoreWeave capped at \$6.3B running to 2032 and a 9.3\% passive stake in Nebius; and, on reporting rather than on a filing, a \$1.5B arrangement with Lambda and a revenue-sharing programme for neoclouds~\cite{nvda-backstop} [A]. And on 10~August 2026 NVIDIA announced memoranda with Apollo, BlackRock, Blackstone, Brookfield, Goldman Sachs and KKR to establish financing platforms mobilizing ``over \$500 billion of third-party capital''~\cite{nvda-500b} [P]---memoranda, it should be said, expressly subject to execution of final agreements.

The regulatory precondition is worth naming because it is almost never reported alongside the announcement. On 29~July 2026, twelve days earlier, the SEC's Office of Structured Finance took the position that securities in datacentre securitizations are not asset-backed securities under the Exchange Act, because datacentres are ``tangible, physical assets that endure beyond the tenor of the securities'' rather than self-liquidating financial assets---placing them outside Regulation~AB and outside the five per cent credit-risk-retention requirement~\cite{sec-dcs} [V]. The financing structure became cheap to build a fortnight before it was built.

NVIDIA's own framing of what this is for could hardly be more explicit. Huang: the company ``began by building chips; today, we are helping create a new class of productive, investable infrastructure: AI factories,'' and NVIDIA compute is ``fungible and transferable across customers and operators.'' Goldman Sachs' chief executive, on the same release: ``a market for credit backed by NVIDIA compute''~\cite{nvda-500b}. On the earnings call: ``Independent capital still underwrites every deal on its own merits. \emph{We're not making loans.} In this model, we get paid twice---once on the hardware sale, and again through the share of rental revenue''~\cite{nvda-defence} [P].

Read the two moves together and the shape is unmistakable. If the model layer is a commons that anyone can serve, and the datacentre layer is financed by capital markets against the residual value of the accelerator rather than against the creditworthiness of any particular tenant, then the failure of any individual model company does not remove demand for the silicon. The commons survives its sponsors. Firms would appear to serve it in the way that distributions, support vendors and integrators appeared around Linux once the kernel was free and no longer owned by anyone whose bankruptcy could take it away. The qualification introduced above applies here with force, and the reader should carry it forward: the analogy holds only for as long as the commons is genuinely unowned, and the acquisitions of August 2026 are an attempt to own not its licence---which cannot be owned---but its means of production. That is a coherent strategy for a company whose exposure is not to which laboratory wins but to whether inference happens at all, and it is the mirror image---in method, not in beneficiary---of the argument Chapter~\ref{ch:model} makes about open weights.

\subsection{The counter-evidence, which is substantial}

The book's discipline requires the other half, and here it is unusually strong. Four things cut against reading this as a settled, confident strategy.

First, it was partly suspended. On 27~August 2026, one day after the earnings call, NVIDIA paused parts of the July revenue-sharing and credit-support programme, reportedly after employees flagged potential antitrust exposure and after friction over how far NVIDIA could dictate which customers its partners rented to~\cite{nvda-pause} [A]. NVIDIA's own statement is that the model ``is still in place and continues to evolve due to high demand''~\cite{nvda-pause} [P]. A programme parts of which are halted within eight weeks of launch is not the execution of a long-laid plan.

Second, the trajectory is retrenchment rather than escalation---though the most-repeated version of that claim is half the story. A \$100B commitment to OpenAI announced in September 2025 became a \$30B equity position; talks reported at up to \$250B for the Ohio structure became a reported figure under \$120B covering only the first phase, and then \$105B on the filing. The half usually dropped is that while the equity commitment fell by seventy per cent, the \emph{compute} commitment rose: on the second-quarter call Huang put OpenAI's ``existing and planned commitments'' at approximately twelve gigawatts of NVIDIA compute through 2030~\cite{nvda-huang-q2} [P]. NVIDIA reduced what it was prepared to own of its largest customer and increased what it expects to sell that customer. Whether that reads as retrenchment or as discipline depends on which of the two numbers one thought was the strategy, and the popular framing---``NVIDIA cut its OpenAI bet by seventy per cent''---reports only the number that fell. Third, the structures themselves are narrower than the headline. The Ohio guarantee \emph{terminates early if OpenAI achieves a satisfactory credit rating}, and OpenAI must reimburse and indemnify NVIDIA for every amount paid~\cite{nvda-ohio}---which makes it, structurally, a bridge to creditworthiness rather than permanent underwriting. Management denies the vendor-financing characterization directly, and has done since 2025: ``unlike Lucent, NVIDIA does not rely on vendor financing arrangements to grow revenue''~\cite{nvda-critique}.

Fourth, and against all of the above, there is one number in the filings that supports the circular reading better than any critic's argument. In the second quarter of fiscal 2027, NVIDIA's GAAP diluted earnings per share (\$2.46) \emph{exceeded} its non-GAAP figure (\$2.22)---an inversion of the usual relationship---because GAAP other income of \$7.773B was almost entirely gains on equity securities~\cite{nvda-q2fy27} [P]. That pre-tax figure equals roughly an eighth of the quarter's GAAP net income---the after-tax contribution is smaller, since the gain is pre-tax and the income is not---and NVIDIA's disclosed equity portfolio of some \$95.6B is concentrated in the firms that buy its chips~\cite{nvda-exposure}. That is not an accusation; it is a line in a filing, and it is the single most quantitatively honest statement of the concern available. The best-argued external version is Vellante's: independent capital can fund more factories, but independent capital is not independent demand, and the very need for credit enhancement is evidence that the market does not yet accept compute as enduring collateral unaided~\cite{nvda-critique} [A]. Morgan Stanley initiated credit coverage at Neutral, estimating roughly \$200B of total credit exposure by end-2028 of which about \$170B off balance sheet, and coined ``balance-sheet-as-a-service''~\cite{nvda-morgan-stanley} [A].

\subsection{The quarter, as the filings state it rather than as it was reported}

One correction of the record belongs here, because the most closely watched quarter in the sector was reported with its most informative disclosure missing. NVIDIA's second-quarter fiscal 2027 CFO commentary changed the reporting structure, and almost all coverage passed over it~\cite{nvda-segments} [P]. ``Gaming'' has ceased to be a reported market platform; GeForce is folded into a new \emph{Edge Computing} line at \$7{,}198M, up twenty-seven per cent. More consequentially, Data Center is now reported in two pieces: \emph{Hyperscale} at \$48{,}710M, up 102 per cent, and \emph{AI Clouds, Industrial \& Enterprise} at \$40{,}313M, up 138 per cent, with management stating that the second ``will represent roughly half of our data center business'' going forward. The split is itself evidence, and it bears directly on Section~\ref{sec:mirror-silicon}: the faster-growing half of NVIDIA's datacentre business is \emph{not} the hyperscalers who are building their own accelerators. If custom silicon is a substitution threat, it is a threat concentrated in the slower-growing half of a business whose other half compounds at 138 per cent---which does not make the threat unreal, but does locate it.

Three further items are worth stating precisely because they are usually stated loosely. Networking was described on the call as a record quarter growing eighteen per cent sequentially, and \emph{no dollar figure was published}: a company that volunteered a datacentre split it was under no obligation to volunteer withdrew the number for the layer this chapter argues is being substituted first, and the asymmetry deserves recording rather than interpreting. Hopper shipments to China were under one per cent of datacentre revenue, and the outlook assumes \emph{no} China datacentre compute revenue at all---the cleanest available statement both of what Chapter~\ref{ch:trade}'s controls have accomplished and of how little of NVIDIA's forward number now depends on their relaxation. And management guided on the call to roughly seventy per cent revenue growth in fiscal 2028 against a sell-side consensus near forty-four~\cite{nvda-segments} [P]---a forward statement that appears in the transcript and not in the press release, and on any reading the most consequential thing said in the quarter. It is guidance and not a result, offered here as what management believes rather than as what will occur [S]; a company that expects seventy per cent growth two years out is not a company that believes it is being commoditized, and a company that guides to it publicly has also made itself falsifiable, which is more than most of the claims in this chapter can say.

Huang's own framing of the demand side names the thing this chapter is about. Listing the drivers of the buildout, he sets ``a thriving open-model ecosystem'' alongside multiple frontier laboratories scaling in parallel and physical AI coming online, and compresses the entire argument of Section~\ref{sec:mirror-move} into three words: ``compute is revenue''~\cite{nvda-huang-q2} [P]. A chip vendor that lists the open-model ecosystem among its demand drivers has explained, without being asked, why it would buy one.

\subsection{AMD runs the same playbook, which matters for the diagnosis}

AMD releases Instella and Instella-MoE as \emph{fully} open models---weights, data and training code---positions ROCm~7 explicitly against CUDA, and chairs the UALink consortium whose promoter group is every major buyer of accelerators and does not include NVIDIA~\cite{amd-open}. On the financing side it has issued warrants for up to 160 million shares at one cent to OpenAI, and an equivalent 160-million-share warrant to Meta, vesting against gigawatt purchase tranches and escalating share-price thresholds; committed up to \$5B of equity to Anthropic against 2\,GW of deployment; and disclosed \$4.1B of datacentre lease guarantees~\cite{amd-warrants,amd-q2-2026} [P]. Stated plainly, the customer is being offered payment in the vendor's own equity for buying the vendor's product---though as of the August 2026 filing no warrant shares had vested or been issued~\cite{amd-warrants}.

That AMD runs the same playbook is analytically important, because it argues for competitive necessity rather than grand design. When two rivals independently adopt an unusual instrument, the more parsimonious explanation is that the instrument is what the situation requires, not that both had the same idea about the future shape of the industry. The book should hold both readings; Section~\ref{sec:mirror-license} says what each licenses.

\section{The laboratories and the hyperscalers commoditize the silicon}
\label{sec:mirror-silicon}

The move in the other direction is equally overt, and Hot Chips 2026 was where it became impossible to describe as a hedge~\cite{hc2026}. Appendix~\ref{app:hotchips} carries the ledger; the strategic content is this.

OpenAI presented Jalape\~no, an inference accelerator co-designed with Broadcom: 13.4 PFLOP/s at four bits, 216\,GiB of HBM4 at 15.4\,TB/s, a 128-chip scale-up domain and a 2{,}048-chip pod, under a term sheet covering ten gigawatts of deployment through 2029~\cite{hc2026-jalapeno,openai-broadcom} [P]. Google announced, for the first time, a training part and an inference part developed concurrently rather than alternately, and---the more consequential fact---that it would \emph{begin} delivering TPUs to a select group of customers in those customers' own datacentres, with a small share of that hardware revenue in 2026 and the majority expected in 2027~\cite{hc2026-tpu8,tpu-external} [P]. Microsoft's Maia~200 makes an economic rather than an architectural claim: thirty per cent better performance per dollar than the latest generation in its fleet, with NVIDIA unmentioned~\cite{maia200}. Meta has disclosed a four-generation MTIA roadmap on roughly a six-month cadence---of which one generation is in production, the next in lab testing---with hundreds of thousands of earlier parts already deployed~\cite{mtia400}. AWS states more than a million Trainium processors deployed and calls it a multi-billion-dollar business~\cite{trainium3}. Behind all of them sits Broadcom, whose AI semiconductor revenue grew 143 per cent to \$10.8B in its second fiscal quarter with the next quarter guided above \$16B~\cite{broadcom-q2fy26} [P].

A defensible order of magnitude follows from two primary filings rather than from any analyst's estimate. Annualizing Broadcom's guided AI semiconductor run rate (above \$16.0B a quarter, so above \$64B a year) against NVIDIA's datacentre run rate (\$89.0B in the quarter, so about \$356B a year) puts merchant custom-ASIC revenue at roughly fifteen per cent of the combined total; annualizing against NVIDIA's \emph{total} third-quarter guidance of \$108.0B instead, of which datacentre is the overwhelming majority, puts it nearer fourteen; taking both companies' most recent actuals rather than their guidance puts it nearer eleven. Call it eleven to fifteen per cent, and note that the range is set by which quarter and which basis one chooses rather than by any disagreement about the numbers [M]. Custom silicon is not yet a substitute for the merchant GPU. It is, however, growing several times faster, and it is being bought by precisely the firms whose purchase orders set NVIDIA's revenue.

The fabric tells the same story more sharply, because there the substitution has already happened. Among the parts presented with an off-package scale-up fabric, the Ethernet-class designs are the clear majority---AMD's UALink-over-Ethernet on 512-port merchant switch ASICs, Microsoft's single Ethernet protocol from chip to a 6{,}144-accelerator Clos under a slide titled ``Why Unified?'', SambaNova's ten integrated 800G ports per socket, Meta's RoCE on dedicated I/O chiplets, and OpenAI's 2{,}048-chip pod over a Broadcom Tomahawk~6 Clos. Of the alternatives, Google's optically switched torus is proprietary but is not sold, and NVIDIA's NVLink is the only proprietary scale-up fabric a third party can buy~\cite{hc2026-scaleup,hc2026-jalapeno,hc2026-tpu8}. The layer that was the deepest source of lock-in three years ago is being rebuilt out of parts anyone can buy.

\subsection{Jalape\~no as a direct challenge, and the sentence that constrains how it can be read}

The Hot Chips presentation was an architecture; three weeks later there were numbers. On 25~August 2026 OpenAI published first results, and they are the sharpest claim any custom part has yet made against the merchant GPU: 1.5 to 1.9 times more AI work per watt at peak throughput, and 1.7 to 3.6 times lower end-to-end latency, measured against GB200 and GB300 systems and against AMD's MI355X, on GPT-OSS 120B, DeepSeek R1 670B and Kimi K2.5~\cite{jalapeno-results} [P]. Two qualifications belong in the same breath and are not decoration. The part \emph{was not deployed at the currency date}---OpenAI says it will ``begin deploying'' Jalape\~no within its own compute infrastructure by the end of the year, which means every figure above is a vendor-run measurement of a part in no production service. And SemiAnalysis observes that the comparison sets HBM4 parts against older HBM3E systems, calling it ``somewhat incomplete and unfair''~\cite{jalapeno-results} [A]. ``Beats NVIDIA by 1.9$\times$'' is therefore a contested claim rather than a settled one, and this book carries it as contested. Nor does it carry a figure for how much of OpenAI's inference the part will serve: \emph{OpenAI has not stated one}, the temptation to supply an estimate should be resisted, and the whole commercial significance of the programme turns on the share that has not been stated.

What the same post does state is the sentence that constrains every reading of it: ``We will continue to widely deploy accelerators from NVIDIA and other partners for both training and inference workloads''~\cite{jalapeno-results} [P]. Set it beside Huang's twelve gigawatts of OpenAI commitments through 2030~\cite{nvda-huang-q2} and the two firms are saying the same thing from opposite directions---the challenger says it will keep buying, and the incumbent says it has been told so. That is the shape this chapter has argued for throughout: a layer being commoditized by a firm that goes on buying from it. It is not the shape of displacement, and the distinction between the two is the whole content of Section~\ref{sec:mirror-license}'s second prohibition.

The analyst reading divides cleanly and both halves are worth printing. Alexander Harrowell of Omdia calls Jalape\~no ``the biggest competitive threat to NVIDIA,'' on the ground that roughly half of AI infrastructure spending now comes from firms running or capable of running custom-silicon programmes; Adrien Sanchez of Yole locates the threat more precisely, at NVIDIA's \emph{inference} margins rather than at its training franchise---the more defensible version, and the one this chapter adopts as the live hypothesis, since inference is where a fixed-function part's advantages survive contact with a workload. Against them, Daniel Newman of Futurum calls the ``NVIDIA killer'' framing ``exhausting'' and credits Broadcom rather than OpenAI for the engineering, which has more force than it first appears: if the engineering is Broadcom's, then the rent displaced from the merchant GPU lands at a merchant ASIC vendor rather than at a laboratory, and the laboratory has commoditized its own supplier only to acquire a different one~\cite{jalapeno-analysts} [A]. Huang's answer is the incumbent's classic one and is recorded here rather than dismissed: ``Lots of projects get started. Lots of projects get canceled\ldots\ We have the supply chain and the technology scale to be their largest supplier''~\cite{jalapeno-analysts} [P].

One feature of the episode is easy to miss and matters to a book about strategy. \emph{Neither OpenAI's announcement nor Broadcom's names NVIDIA.} The entire adversarial framing---the comparisons, the ``killer'' language, the threat assessment---is supplied by analysts and by the press. The challengers decline to describe the challenge as one, which is either commercial courtesy toward a supplier from whom they still expect twelve gigawatts, or an accurate account of an intent that was never adversarial; from outside the firms the two are not distinguishable. This chapter records the framing as third-party framing, because that is what it is, and notes that the discipline cuts both ways---a challenge is not less real for going unnamed, and it is not more real for being named by someone with no part in it.

A correction of the record belongs here too, dryly, because a false claim circulates alongside this story and the true version is more useful. It is repeated that NVIDIA's head of datacentre left for OpenAI. He did not, and the movement ran the other way: \emph{OpenAI's} own Head of Data Centers, Chris Malone, left \emph{OpenAI} in the week of 17~August 2026---reported on the 25th---after sixteen months, amid a continuing stream of senior departures. The provenance is worth stating precisely, because it is commonly garbled: the ``head of data centers'' title was OpenAI's, not Meta's---at Meta he spent close to five years as a \emph{distinguished engineer}, and before that nearly twelve years at Google on datacentre infrastructure and IT hardware research---and he had never worked at NVIDIA~\cite{openai-exec} [V]. That career shape is itself the point: the person OpenAI recruited to build its datacentres was an infrastructure engineer of seventeen years' standing at the two firms that had built the largest ones, and OpenAI kept him for sixteen months. OpenAI's infrastructure organization now reports to Sachin Katti, recruited in November 2025 from Intel, where he was chief technology and AI officer. NVIDIA's datacentre leadership is intact---Ian Buck remains VP of Hyperscale and HPC---and the senior technical flow of 2026 ran \emph{toward} NVIDIA: Groq's leadership and roughly ninety per cent of its staff in December 2025, then 109 Poolside researchers in August 2026~\cite{nvda-structured,nvda-poolside}. This is not scorekeeping. A book that takes the custom-silicon challenge as seriously as this chapter does has to get the direction of the flow right, because a labour market moving toward the incumbent is not evidence of the incumbent losing the argument---and it is not evidence of the incumbent winning it either, since NVIDIA paid roughly \$27B for those two flows and a market in which talent must be bought at that price is not obviously a market at rest.

\subsection{And the counter-evidence here is decisive enough to be printed first in any honest summary}

On 26~August 2026, in the last week before this edition's currency date, AWS and NVIDIA announced two million additional NVIDIA GPUs across 2027--2028---and, materially, that \emph{Trainium4 will adopt NVIDIA's NVLink Fusion interconnect and NVIDIA's custom HBM}~\cite{aws-nvda} [P]. The largest deployer of non-NVIDIA AI silicon in the world is placing its own accelerator inside NVIDIA's rack-scale interconnect and memory stack. Set against NVIDIA's quarter---datacentre revenue up 117 per cent to \$89.0B, against total-company third-quarter guidance of \$108.0B $\pm$2 per cent~\cite{nvda-q2fy27}---the honest reading is not that the incumbent is being displaced by the custom-silicon wave but that it is absorbing it. The commoditization of the silicon layer is real, it is well funded, and it has so far moved the boundary of what NVIDIA sells rather than the size of what NVIDIA sells.

\section{And one actor commoditizes the fab}
\label{sec:mirror-fab}

At the bottom of the stack, SpaceX confirmed in August 2026 a \$16.8B initial investment in ``Terafab,'' a vertically integrated logic and memory manufacturing, packaging and test site in Grimes County, Texas, with output targeted at Tesla's robotics and vehicle programmes and at space-based datacentres~\cite{terafab}. Separately, and not disclosed as part of Terafab, Intel entered a pact with Tesla, SpaceX and xAI in April 2026 covering advanced logic and packaging, which one analyst reads as an Intel fab expansion with those three as anchor customers rather than a standalone venture~\cite{terafab} [A]. In parallel Tesla runs a dual-foundry arrangement for AI5 and AI6 at TSMC and Samsung, with the Samsung Taylor 2\,nm fab reaching AI5 tape-out in July 2026~\cite{tesla-ai5}.

This one deserves the coldest grade in the chapter. The announced figures have moved three times---from a \$20--25B analyst estimate in March, to \$55B initial and \$119B total in a May regulatory filing, to \$16.8B initial in the August confirmation. The process node was not restated in the confirmation, no construction or production timeline was disclosed, and the project was not discussed on SpaceX's first earnings call in the same week. Terafab belongs in this chapter as evidence of intent [R], not as capacity. What it demonstrates is the reach of the logic: an actor whose rent is in launch, vehicles and robots, facing a fabrication layer it regards as a constraint, reasons its way to owning the constraint. That is move five of the playbook---supply-chain capture---executed by an American company against an American and Taiwanese supply chain.

\section{The same method, and the asymmetry that survives it}
\label{sec:mirror-asymmetry}

Table~\ref{tab:mirror} sets the actors side by side. The columns that matter are the second and third: what each is defending, and what each is willing to burn to defend it.

\begin{table}[htbp]
\centering\footnotesize
\caption{Who commoditizes what, as of \datafreeze. ``Instrument'' lists the specific vehicle, not the intent. ``Commoditized'' is used broadly, and the breadth matters: it covers both \emph{giving a layer away}, which is what NVIDIA and AMD do at the model layer, and \emph{substituting for it in-house}, which is what a hyperscaler building its own accelerator does---the first destroys a rent for everyone, the second only removes the actor from paying it. The two are different instruments with the same effect on the layer's pricing power, and the last two rows are of the second kind. Status: \textbf{running} (executing and observable), \textbf{partial} (executing with a component suspended or unexecuted), \textbf{announced} (committed, not yet productive).}
\label{tab:mirror}
\begin{tabular}{@{}p{25mm}p{33mm}p{33mm}p{40mm}p{15mm}@{}}
\toprule
\textbf{Actor} & \textbf{Layer defended} & \textbf{Layer commoditized} & \textbf{Instrument} & \textbf{Status} \\
\midrule
China (state and firms) & Hardware volume, domestic cloud, standards, dependence & \emph{All of them} & Open weights, open harness, open ISA, open interconnect; mandate; power tariff & Running \\
NVIDIA & Accelerator and CUDA & The model, and the financing of the datacentre & Nemotron weights \emph{and} training data; OpenMDW; cuTile (Apache~2.0); Poolside ``Model Factory'' licence and 109 researchers; reported Hugging Face acquisition [A]; residual value guarantees; revenue share & Partial (revenue share paused Aug.) \\
AMD & Accelerator & The model; the scale-up fabric & Instella (weights, data, code); ROCm; UALink over merchant Ethernet; customer warrants & Running \\
OpenAI & The frontier model and the platform & The accelerator and the system & Jalape\~no with Broadcom; 10\,GW term sheet & Running \\
Google, Microsoft, Meta, AWS & Cloud platform and applications & The accelerator, the system, the fabric & TPU 8t/8i sold into customer sites; Maia~200; MTIA; Trainium; Ethernet scale-up & Running \\
Broadcom, Marvell & Co-design and networking & The merchant GPU & Custom ASIC programmes; eRoCE & Running \\
SpaceX, Tesla & Launch, vehicles, robots & The fab & Terafab; dual-foundry AI5/AI6 & Announced \\
\bottomrule
\end{tabular}
\end{table}

Two asymmetries survive the convergence of method, and they run in opposite directions.

\textbf{The first favours the West, and the book should say so plainly.} Commoditization is not the same as losing. A stack in which every layer is cheap is a stack in which the surplus accrues to whoever holds the layer that stays scarce---and at the currency date the scarce layers are, for the most part, still Western-held: the CUDA installed base, the enterprise deployment surface, the trust and certification apparatus of Chapter~\ref{ch:trust}, and the deployment-generated data of Chapter~\ref{ch:data}. The clearest single demonstration that the West is executing rather than retreating is the price inversion of Section~\ref{sec:price-inversion}: on 30~July 2026 the incumbent cut the bottom tier of its ladder by eighty per cent to \$0.20 and \$1.20 per million tokens, which is \emph{below the list price of the frontier-scale Chinese open-weight models}---though not of their distilled tiers, a restriction Section~\ref{sec:price-inversion} states precisely---and usage of that tier rose to almost fourteen times its pre-cut level with roughly three-quarters of the gain taken from competitors~\cite{openai-luna,price-inversion} [P/A]. A firm that can price underneath the commodity and gain share doing it is not a firm being commoditized. It is a firm commoditizing.

\textbf{The second favours China, and it is structural rather than tactical.} Every Western move in Table~\ref{tab:mirror} commoditizes a layer in which \emph{another Western firm} was earning the rent. NVIDIA is destroying the model layer's rent in order to protect silicon. OpenAI and the hyperscalers are destroying the silicon layer's rent in order to protect the model and the platform. Each decision is individually rational, and the honest statement of what they do jointly has to respect the mechanism of Section~\ref{sec:mirror-move}: they \emph{redistribute} the Western rent pool rather than necessarily shrinking it, because the surplus a commoditizing move destroys at one layer is the surplus it delivers to another. The pincer claim therefore has to be conditional, and this book states it conditionally: it bites to the extent that the layers the rent moves \emph{to} are layers a Chinese actor can also occupy, and not otherwise. What China gains from the arrangement is not the destruction of a rent pool but a cheaper entry condition. It does not need to destroy either layer; it needs only to be present in both when the rents have moved, and presence is very much cheaper than victory.

Sharper still, and this is the sentence the chapter exists to reach: \emph{the layers the West is commoditizing are disproportionately the layers where China is behind, and several of the commoditizing moves directly reduce the value of Western state instruments.} Open weights released with their training data lower the cost of closing the model-layer gap for whoever is closing it. An open instruction set that NVIDIA has publicly qualified as a CUDA host lowers the cost of the compute-layer detour that Chapter~\ref{ch:compute} documents~\cite{hc2026-nv-riscv}. A scale-up domain rebuilt out of merchant Ethernet switch silicon is precisely the substitution that the chokepoint strategy was designed to prevent, and it is being performed by the firms the chokepoint strategy was designed to protect~\cite{hc2026-scaleup}. None of these is a concession to China; all of them are commercial decisions taken against a domestic rival. The externality is the same either way. A denial policy that assumes the denied party must obtain the controlled item from the controlling country is undermined as effectively by an American open-weight release as by a Chinese one, and the American release comes with the training data attached.

\subsection{What this does to Proposition~P3}

The book's transmission argument for P3 required Chinese price pressure to reach Western margins: cheap open weights would force the closed frontier to cut, the cut would compress the rent, and the compressed rent would fail to finance the next training run. Section~\ref{sec:price-inversion} shows that chain is weaker than it looked---the incumbent cut its cheapest tier, not its flagship, held the flagship for three weeks, made the flagship cut temporary, and the price leader on the Chinese side \emph{raised} prices by as much as 1{,}100 per cent mid-cascade~\cite{deepseek-raise}. But P3 does not depend on that chain any more. It is now overdetermined in \emph{mechanism}, because the largest supplier of the frontier's inputs is itself releasing open weights with training data at zero price, and the frontier's largest customers are building the silicon underneath it. One qualification is owed, and it is the weakest joint in this chapter's argument, so it is stated here rather than left for a reader to find. NVIDIA's \emph{shipped} open models take routed volume but not frontier capability---the best of them scores well short of the leading open Chinese weights on the independent index (38 against 60 on the August v4.1.1 board, Appendix~\ref{app:bulletin})---so what NVIDIA has demonstrated to date is that a chip vendor can commoditize the \emph{volume} tier of the model layer, not that it can commoditize the capability tier. The stronger claim, that rent destruction at the model layer no longer requires a Chinese actor to arrive, is contingent on a frontier-scale NVIDIA release that at \datafreeze has not happened [S]. The Poolside transaction is the clearest evidence yet that such a release is being funded rather than merely wished for---\$6B for a training platform and 109 researchers is a capability purchase and not a marketing one~\cite{nvda-poolside}---but a funded intention remains an intention [R]. What is not contingent is the direction of travel and the stated intent, and those are on the record.

That is a strengthening of P3's mechanism and simultaneously a weakening of P3's \emph{diagnostic value for this book's thesis}. A proposition that can be satisfied by the incumbent's own conduct is a weaker piece of evidence about who wins, and the register in Appendix~\ref{app:evidence} should be read with that in mind.

\section{What the mirror licenses, and what it forbids}
\label{sec:mirror-license}

It licenses three claims. That the playbook of Part~I is a general theory of a contest over a layered technology stack rather than a description of Chinese exceptionalism---which is a considerable strengthening of Part~I, because a theory that predicts the behaviour of actors it was not built from is doing more work than one that does not. That the commoditization of the AI stack is now overdetermined, proceeding from several directions at once and no longer contingent on any single actor's strategy holding. And that the West's own commercial decisions are eroding the West's own policy instruments, which is a finding Chapter~\ref{ch:strategies}'s instruction to commoditize before being commoditized has to absorb rather than a rhetorical flourish: an export-control regime and an open-weight corporate strategy are not merely in tension, they are aimed at the same object from opposite sides.

It forbids three claims equally. That the West is passively losing---it is not; it is executing hard, it has priced underneath the commodity at the bottom of the ladder and gained share doing so, and its incumbent grew datacentre revenue 117 per cent in the quarter this chapter closes in. That commoditization implies a Chinese victory---it does not; it implies that the rent moves, and whether it moves to a layer China holds is the open question of the whole book rather than a corollary of this chapter. And that NVIDIA's financing structures are straightforwardly vendor financing---management denies it in terms, the largest structure terminates on a credit-rating event and carries full indemnification, one programme was suspended within eight weeks over antitrust concerns, and the honest verdict on the evidence available at the currency date is that the practice is unusual, large, growing, disclosed, and not yet demonstrated to be either sound or unsound.

One item is neither licensed nor forbidden but held open, and it is named here so that a reader can date it. If the reported acquisition of Hugging Face is confirmed, the description of the open-model layer in Section~\ref{sec:mirror-nvidia} changes in kind rather than in degree, and the Linux analogy this chapter uses twice must be withdrawn rather than qualified: a commons whose warehouse belongs to the vendor of the complement is a distribution channel with a permissive licence attached. If it is not confirmed---and at \datafreeze it is not---the Poolside licence and its 109 researchers stand on their own as the strongest evidence available that the open-model layer now has a principal sponsor, which is the weaker claim and still a large one.

Two tests would settle much of this by the end of 2027, and both belong in the register maintained in Appendix~\ref{app:dashboard}. If merchant-GPU datacentre revenue growth decelerates while custom-ASIC revenue continues at Broadcom's guided trajectory, the silicon commoditization is working and the pincer is real. If instead custom silicon continues to converge onto NVIDIA's interconnect and memory---as Trainium4 has already done~\cite{aws-nvda}---then the wave is being absorbed, the accelerator vendor's giveaway of the model layer was correctly priced, and the layer that stays scarce is the one it has held all along.
